Let $\mu$ be a law on $A^{\Z^d}$ such that every one-site conditional atom,
conditioned on all other sites, is at least $\eta>0$.  We give two
cellular-automaton regimes in which this hypothesis yields a sharp cylinder
estimate at every time: nonconstant scalar affine rules over a finite field,
and finite-alphabet rules permutive at a uniquely exposed neighborhood site.
For every $t\geq0$, finite $S\subset\Z^d$, and $a\in A^S$, the image law
satisfies
\[
  (F^t_*\mu)([a]_S)\geq\eta^{|S|}.
\]
No invariance, independence, stationarity, or mixing assumption is imposed.
The probabilistic step upgrades the one-site hypothesis to a conditional
finite-block bound with an outside-measurable target.  Scalar affine rules
then supply $|S|$ site pivots through row independence in a Laurent
polynomial domain.  In the exposed-corner case, the corner remains permutive
under iteration and translated corner pivots form a triangular bijection on
arbitrary finite shapes.  Equality occurs for a least-likely iid pattern
under the identity rule.  We also give a reversible vector-alphabet
counterexample to the analogous whole-site bound and a surjective
nonpermutive rule showing that surjectivity alone does not provide the
triangular hypothesis.

